% !PRINTED 766
\weberpage{766}
Danach ist die ganze rationale Zahl

\[
C_1\,F(z_1)+C_2\,F(z_2)+\cdots+C_m\,F(z_m)
\]

nach dem Modul $p$ mit

\[
c^{p-1}\left[C_1\,\varphi'(x_1)^p+C_2\,\varphi'(x_2)^p+\cdots+C_m\,\varphi'(x_m)^p\right]
\]

congruent. Da nun $C_1, C_2, \ldots, C_p$ ganze rationale Zahlen sind, % [sic] -- printed $C_p$; the surrounding system has $m$ constants
so ist $C_1^p\equiv C_1, C_2^p\equiv C_2, \ldots \pmod{p}$, und es ergiebt sich durch Anwendung des polynomischen Lehrsatzes:

\[
\begin{aligned}
C_1\,\varphi'(x_1)^p+C_2\,\varphi'(x_2)^p+\cdots+C_m\,\varphi'(x_m)^p &\equiv \\
\left[C_1\,\varphi'(x_1)+C_2\,\varphi'(x_2)+\cdots+C_m\,\varphi'(x_m)\right]^p &\equiv k^p \pmod{p}.
\end{aligned}
\]

Danach erhalten wir also die Congruenz

\weq{18}{C_1\,F(z_1)+C_2\,F(z_2)+\cdots+C_m\,F(z_m) \equiv c^{p-1}\,k^p \pmod{p}.}

Nun sind die Zahlen $c, k$ von $p$ unabhängig, und wir können daher $p$ so gross annehmen, dass es nicht in $c$ und in $k$ aufgeht.

5. \gesperrt{Dann ist die Summe $C_1\,F(z_1)+C_2\,F(z_2)+\cdots+C_m\,F(z_m)$ eine von Null verschiedene ganze rationale Zahl, also dem absoluten Werthe nach mindestens gleich 1.}

Wir bezeichnen nun mit $\varphi_1(x)$ die ganze Function, die sich aus $\varphi(x)$ ergiebt, wenn die negativen unter den Coëfficienten $a, a_1, a_2, \ldots$ durch ihre positiven Werthe ersetzt werden. Die Function

\[
f_1(z) = \dfrac{\varphi_1(x)^{p-1}\,\varphi_1'(x)}{\Pi(p-1)},
\]

die sich nach (17) ergiebt, hat dann nur positive Coëfficienten, und diese Coëfficienten sind dem absoluten Werthe nach gewiss nicht kleiner, als die entsprechenden Coëfficienten von $f(z)$, weil die Coëfficienten von $f_1(z)$ aus denselben Zahlen durch Addition entstehen, die bei der Bildung der Coëfficienten von $f(z)$ theils addirt, theils subtrahirt werden.

Hieraus aber ergiebt sich nach der Formel \S.\,205, (10), wenn wir mit $\xi$ den absoluten Werth von $x$ bezeichnen, für den absoluten Werth von $Q(z)$:

\[
|Q(z)| \leqq \dfrac{\varphi_1(\xi)^{p-1}\,\varphi_1'(\xi)}{\Pi(p-1)},
\]

6. \gesperrt{und es kann also $Q(z)$ für jedes endliche $z$ durch hinlängliche Vergrösserung von $p$ beliebig klein gemacht werden.}
